% Partie : intuitionistic-logic
% Chapitre : soundness-completeness
% Section : canonical-model

\documentclass[../../../include/open-logic-section]{subfiles}

\begin{document}

\olfileid[fr]{int}{sc}{mod}

\olsection{Le modèle canonique}

Les mondes de notre modèle seront des suites finies~$\sigma$
d'entiers naturels, c'est-à-dire des éléments
$\sigma \in \Nat^*$. Rappelons que $\Nat^*$ se définit par induction :
\begin{enumerate}
\item $\emptyseq \in \Nat^*$.
\item Si $\sigma \in \Nat^*$ et $n \in \Nat$, alors
  $\sigma.n \in \Nat^*$, où $\sigma.n$ désigne
  $\sigma \concat \tuple{n}$ et $\sigma \concat \sigma'$ la
  concaténation de $\sigma$ avec $\sigma'$.
\item Rien d'autre n'appartient à $\Nat^*$.
\end{enumerate}
Nous pouvons donc utiliser $\Nat^*$ pour donner des définitions par
induction.

Soit $\tuple{!B_0, !C_0}$, $\tuple{!B_1, !C_1}$, \dots, une
énumération de tous les couples de !!{formula}s. Étant donné un
ensemble de !!{formula}s~$\Delta$, définissons $\Delta(\sigma)$
par induction :
\begin{enumerate}
\item $\Delta(\emptyseq) = \Delta$
\item $\Delta(\sigma.n) = {}$
  \[
  \begin{cases}
    (\Delta(\sigma) \cup \{!B_n\})^* &
    \text{si $\Delta(\sigma) \cup \{!B_n\} \Proves/ !C_n$} \\
    \Delta(\sigma) & \text{sinon}
  \end{cases}
  \]
\end{enumerate}
Par $(\Delta(\sigma) \cup \{!B_n\})^*$, nous entendons un ensemble
premier de !!{formula}s choisi parmi ceux dont l'existence est garantie par le
\olref[lin]{lem:lindenbaum}, appliqué à l'ensemble
$\Delta(\sigma) \cup \{!B_n\}$ et à la !!{formula}~$!C_n$.
Par cette définition, si
$\Delta(\sigma) \cup \{!B_n\} \Proves/ !C_n$, alors
$\Delta(\sigma.n) \Proves !B_n$ et
$\Delta(\sigma.n) \Proves/ !C_n$. De plus,
$\Delta(\sigma) \subseteq \Delta(\sigma.n)$ pour tout~$n$.
Si $\Delta$ est premier, $\Delta(\sigma)$ l'est pour tout~$\sigma$.

\begin{defn}\ollabel{defn:canonical-model}
  Supposons $\Delta$ premier. Le \emph{modèle canonique}
  $\mModel{M(\Delta)}$ associé à $\Delta$ est défini ainsi :
  \begin{enumerate}
  \item $W = \Nat^*$ est l'ensemble des suites finies d'entiers
    naturels.
  \item $R$ est l'ordre partiel tel que $R\sigma\sigma'$ si et
    seulement si $\sigma$ est un segment initial de~$\sigma'$,
    c'est-à-dire si $\sigma' = \sigma \concat \sigma''$ pour une
    certaine suite~$\sigma''$.
  \item $V(p) = \Setabs{\sigma}{p \in \Delta(\sigma)}$.
  \end{enumerate}
\end{defn}

Il est facile de vérifier que $R$ est bien un ordre partiel. La
condition de monotonie sur~$V$ est également satisfaite. De
$\Delta(\sigma) \subseteq \Delta(\sigma.n)$, on déduit
$\Delta(\sigma) \subseteq \Delta(\sigma')$ lorsque
$R\sigma\sigma'$ : il suffit de raisonner par récurrence sur la
longueur du suffixe, en partant de~$\sigma$.

\end{document}
